\chapter{The Geometry of Failure: Dimensional Collapse in Low-Bit AI Silicon}

\section{Why High Dimensions Matter}

Modern large language models operate in representation spaces with thousands of
dimensions ($H \ge 1024$).
Unlike physical 3D space, high-dimensional geometry obeys radically different
statistical laws.
These geometric laws fundamentally shape training dynamics, conditioning,
and hardware behavior.

Understanding failure modes in LLM training therefore requires understanding
how high-dimensional random vectors behave under nonlinear transformation,
residual accumulation, and quantized arithmetic.

We begin from first principles.

\section{The Statistical Baseline: Noise Floor of the Mean}

Let $x_i \in \mathbb{R}^H$ be $N$ independent zero-mean random vectors with
finite second moment:

\[
\mathbb{E}[x_i] = 0, 
\qquad
\mathbb{E}\|x_i\|^2 = \sigma^2 H.
\]

The empirical mean is:

\[
\mu = \frac{1}{N} \sum_{i=1}^N x_i.
\]

Even if the population mean is exactly zero, finite sampling introduces
a random offset. By concentration of measure,

\[
\mathbb{E}\|\mu\|^2 = \frac{\mathbb{E}\|x_i\|^2}{N}.
\]

Define the \textbf{noise floor scale}:

\[
\|\mu\|_{\text{noise}} 
\sim \sqrt{\frac{\mathbb{E}\|x_i\|^2}{N}}.
\]

We define the dimensionless diagnostic:

\[
R = \frac{\|\mu\|}{\sqrt{\mathbb{E}\|x\|^2 / N}}.
\]

Interpretation:

\begin{itemize}
\item $R \approx 1$ : mean is consistent with random fluctuation.
\item $R \gg 1$ : mean is statistically structured (coherent bias).
\end{itemize}

This metric allows us to distinguish genuine geometric drift from
finite-sample noise.

\section{The Counter-Intuitive Geometry of High Dimensions}

High-dimensional space ($H \gg 1$) behaves fundamentally differently from
three-dimensional intuition. Two geometric principles are particularly
important for understanding LLM training dynamics.

\subsection{The Orthogonality Paradox}

Let $u, v \in \mathbb{R}^H$ be independent isotropic random vectors.
Then

\[
\mathbb{E}[\langle u, v \rangle] = 0,
\qquad
\mathrm{Var}(\langle u, v \rangle) = O\!\left(\frac{1}{H}\right).
\]

As $H$ increases, random vectors become almost surely orthogonal.
In 1,024 dimensions, two randomly drawn directions are extremely unlikely
to align by chance.

This geometric ``blessing of orthogonality'' allows large language models
to represent many independent semantic concepts simultaneously.
Distinct features (e.g., \textit{apple}, \textit{gravity}, \textit{democracy})
can coexist without interfering because random semantic directions naturally
separate.

\subsection{The Thin-Shell Phenomenon}

For isotropic vectors $x \in \mathbb{R}^H$ with finite second moment,

\[
\|x\| \approx \sqrt{\mathbb{E}\|x\|^2}
\]

with vanishing relative variance as $H$ grows.

Thus, almost all mass concentrates near a hyperspherical shell.
The interior volume becomes negligible.
In high dimension, there is effectively no ``middle'' —
almost everything lies on the surface.

\subsection{Dimensionality as a Magnifier}

Because random alignment variance scales as $1/\sqrt{H}$,
even a small coherent bias becomes statistically amplified.
Define the noise-normalized mean magnitude

\[
R = \frac{\|\mu\|}{\sqrt{\mathbb{E}\|x\|^2/N}}.
\]

In high dimension:

\begin{itemize}
\item $R \approx 1$ corresponds to pure finite-sample noise.
\item $R = 1.1$ is already statistically significant.
\item $R \gg 1$ indicates structural drift.
\end{itemize}

With $H=4096$, accidental alignment across thousands of coordinates is
astronomically unlikely.
Thus even small deviations represent real geometric structure.

\section{Nonlinear Operators as Bias Generators}

Consider zero-mean input $x$.

\subsection{ReLU}

For symmetric $x$,

\[
\mathbb{E}[\text{ReLU}(x)] > 0.
\]

ReLU removes negative half-space, introducing positive drift.

\subsection{GELU / SwiGLU}

Smooth nonlinearities preserve more symmetry but still produce

\[
\mathbb{E}[\phi(x)] \neq 0.
\]

Empirically:

\[
R_{\text{post-activation}} \gg 1
\]

even when input $R \approx 1$.

\subsection{Softmax Attention}

Softmax enforces positivity:

\[
\alpha_i = \frac{e^{z_i}}{\sum_j e^{z_j}}.
\]

All outputs lie in the simplex.
Thus attention layers introduce coherent bias directions.

\section{Residual Accumulation and Mean Drift}

Consider residual recursion:

\[
X_{\ell+1} = X_\ell + F(X_\ell).
\]

If $F$ produces mean drift $\Delta \mu_\ell$:

\[
\mu_{\ell+1} = \mu_\ell + \Delta \mu_\ell.
\]

Therefore:

\[
\mu_L = \sum_{\ell=0}^{L-1} \Delta \mu_\ell.
\]

Unless explicitly removed, bias accumulates linearly in depth.

This produces geometric collapse:

\[
X_\ell \approx \mu_\ell + \text{isotropic noise}.
\]

\section{Dimensional Collapse: The Needle in the Sphere}

When the coherent mean component grows such that $R \gg 1$,
the representation undergoes \emph{dimensional collapse}.

\subsection{Loss of Volume}

Let the activation decompose as

\[
x = \mu + \varepsilon,
\]

where $\varepsilon$ is approximately isotropic.

Total variance satisfies

\[
\mathbb{E}\|x\|^2 = \|\mu\|^2 + \mathbb{E}\|\varepsilon\|^2.
\]

If $\|\mu\|^2$ dominates,
the representation no longer spreads evenly across all $H$ dimensions.
Instead, most magnitude lies along a single dominant direction.

\subsection{Needle vs Sphere}

A well-conditioned representation resembles a hypersphere:
variance is distributed across many orthogonal directions.

Under strong bias, the geometry becomes highly elongated:
vectors concentrate inside a narrow cone aligned with $\mu$.
The effective dimensionality collapses.

Geometrically:
\begin{itemize}
\item Healthy state: nearly isotropic sphere.
\item Collapsed state: elongated needle.
\end{itemize}

\subsection{Feature Erasure}

Suppose semantic distinctions between concepts
(e.g., \textit{apple}, \textit{orange}, \textit{banana})
exist within the isotropic component $\varepsilon$.

If $\|\mu\| \gg \|\varepsilon\|$,
then elementwise magnitude is dominated by $\mu$.
Relative semantic differences become small perturbations
around a large shared offset.

Quantitatively,
the signal-to-bias ratio scales as

\[
\frac{\|\varepsilon\|}{\|\mu\|}.
\]

As this ratio shrinks,
fine-grained distinctions become numerically invisible.
In low-bit arithmetic, these distinctions fall below quantization
resolution, leading to degraded semantic fidelity and hallucination.

Dimensional collapse is therefore not merely aesthetic.
It directly limits representational capacity.

\section{Low-Bit Hardware Amplification of Bias}
\label{sec:lowbit_amplification}

In this section we analyze, from first principles, why mean bias and spectral anisotropy become especially problematic under low-bit arithmetic.
The key message is simple but subtle:

\begin{quote}
High dimensionality amplifies small numerical errors, and when dynamic range is dominated by a coherent mean component, quantization precision is largely wasted tracking bias instead of semantic detail.
\end{quote}

We develop this argument step by step.

%------------------------------------------------------------
\subsection{Quantization Error in High Dimensions}

Consider a vector $x \in \mathbb{R}^H$ that is quantized elementwise.
Let the quantized value be
\[
\hat{x}_i = x_i + \varepsilon_i,
\]
where $\varepsilon_i$ denotes quantization error.

We assume:
\begin{itemize}
\item $\mathbb{E}[\varepsilon_i] = 0$ (unbiased rounding),
\item $\mathbb{E}[\varepsilon_i^2] = \sigma_\varepsilon^2$,
\item errors are approximately independent across coordinates.
\end{itemize}

The total squared perturbation is

\[
\|\varepsilon\|^2 = \sum_{i=1}^H \varepsilon_i^2.
\]

Taking expectation:

\[
\mathbb{E}\|\varepsilon\|^2
=
\sum_{i=1}^H \mathbb{E}[\varepsilon_i^2]
=
H \sigma_\varepsilon^2.
\]

Therefore,

\[
\|\varepsilon\| \sim \sqrt{H} \, \sigma_\varepsilon.
\]

\paragraph{Interpretation.}

Even if each coordinate error is tiny, the total Euclidean shift scales as $\sqrt{H}$.
For $H = 4096$,
\[
\sqrt{4096} = 64.
\]
Thus a seemingly small per-coordinate rounding error is geometrically magnified.

This is a purely high-dimensional phenomenon and does not depend on model architecture.

%------------------------------------------------------------
\subsection{Bias-Dominated Dynamic Range}

Now consider the decomposition of an activation vector:

\[
x = \mu + s + \eta,
\]

where:
\begin{itemize}
\item $\mu$ = coherent mean bias (rank-1 component),
\item $s$ = low-rank spike subspace,
\item $\eta$ = high-dimensional semantic tail.
\end{itemize}

Empirically we observe:
\[
\|\mu\| \gg \|\eta\|.
\]

This implies that elementwise magnitudes are dominated by bias:
\[
|x_i| \approx |\mu_i|.
\]

%------------------------------------------------------------
\subsection{Blockwise Quantization and Scale Selection}

Modern LLM quantization uses blockwise scaling.
For each block $b$ of coordinates, scale is typically

\[
\alpha_b = \max_{i \in b} |x_i|.
\]

If bias dominates,

\[
\alpha_b \approx \max |\mu_i|.
\]

This means:
\begin{itemize}
\item Quantization scale aligns with bias magnitude.
\item Semantic variation is measured relative to this large scale.
\end{itemize}

Thus the effective resolution available for $\eta$ becomes very small.

%------------------------------------------------------------
\subsection{FP4 Precision Allocation}

FP4 provides only 16 representable values.

If scale is $\alpha$, then bin width is approximately

\[
\Delta \approx \frac{2\alpha}{16}.
\]

If semantic variation satisfies

\[
|\eta_i| \ll \alpha,
\]

then semantic coordinates occupy only a few quantization bins.

\paragraph{Bit-budget distortion.}

Most representable levels are used to track the bias component.
Very few are left for semantic structure.

This leads to:
\begin{itemize}
\item angular distortion,
\item loss of discriminative resolution,
\item effective feature erasure.
\end{itemize}

%------------------------------------------------------------
\subsection{Dimensional Collapse}

Consider the geometric picture.

Healthy representation:
\[
x \approx \text{sphere in } \mathbb{R}^H.
\]

Bias-dominated representation:
\[
x \approx \mu + \eta,
\quad \|\mu\| \gg \|\eta\|.
\]

Vectors cluster near direction $\mu$.
The cloud becomes a narrow cone.

After quantization:
\[
\hat{x} = x + \varepsilon.
\]

If $\|\eta\|$ is small relative to $\alpha$,
then $\varepsilon$ may dominate $\eta$.

Geometrically:
\begin{itemize}
\item Sphere $\rightarrow$ Cone.
\item Cone $\rightarrow$ Needle.
\end{itemize}

We call this \emph{Dimensional Collapse}.

When collapse occurs:
\begin{itemize}
\item Representations lose angular diversity.
\item Concepts become indistinguishable.
\item Models hallucinate or degrade.
\end{itemize}

%------------------------------------------------------------
\subsection{Why Mean Removal Fixes Most of the Problem}

If we subtract mean:

\[
x' = x - \mu,
\]

then:
\begin{itemize}
\item Dynamic range shrinks.
\item Quantization scale aligns with spike or semantic magnitude.
\item FP4 bins now allocate precision to meaningful variation.
\end{itemize}

Quantization noise now scales with semantic energy rather than bias magnitude.

%------------------------------------------------------------
\subsection{Why Spike Whitening Helps Further}

After mean removal:

\[
x' = s + \eta.
\]

If spike singular values still dominate,
semantic components still receive less resolution.

Whitening rescales:

\[
x'' = U \Sigma^{-1} U^\top x'.
\]

This flattens the spectrum while preserving information.

Unlike removal, whitening does not destroy structure.

%------------------------------------------------------------
\subsection{Summary}

Low-bit hardware amplifies bias through three mechanisms:

\begin{enumerate}
\item High dimensionality magnifies rounding error by $\sqrt{H}$.
\item Bias dominates quantization scale.
\item Blockwise quantization spreads bias globally.
\end{enumerate}

Mitigation:

\begin{itemize}
\item Mean subtraction removes dominant distortion.
\item Spike whitening restores spectral balance.
\item Both are hardware-efficient and localizable at embedding stage.
\end{itemize}

This provides a geometric explanation for why
bias mitigation dramatically improves low-bit training stability.


\section{Operator-Level Stabilization Principles}
\label{sec:operator_stabilization}

The empirical and theoretical analysis above establishes that
coherent mean bias is regenerated by nonlinear operators
and accumulates through residual pathways.
In low-bit regimes, this bias becomes the dominant source of
numerical distortion.

We now derive operator-level stabilization mechanisms
that directly follow from high-dimensional geometry,
mean propagation dynamics, and hardware arithmetic constraints.

% ============================================================
\subsection{1. Online Centering: Intra-Block Mean Removal}
\label{subsec:online_centering}

Standard Transformer architectures remove mean implicitly at residual
boundaries through LayerNorm.
However, for low-bit arithmetic, this is too late.

The bias that destroys precision is not the residual mean,
but the mean generated \emph{inside} nonlinear operators
before quantization.

\paragraph{Mechanism of Mean Regeneration.}

Let $X \in \mathbb{R}^{N \times H}$ be zero-mean.
After a non-odd activation $\phi$ (ReLU, GELU, SwiGLU),
the output

\[
Y = \phi(X)
\]

satisfies

\[
\mathbb{E}[Y] \neq 0.
\]

Empirically, the noise-normalized ratio

\[
R = \frac{\|\mu_Y\|}{\sqrt{\mathbb{E}\|Y\|^2/N}}
\]

becomes $R \gg 1$.

In FP4 (16 representable values),
this structural shift pushes activations toward
the edges of the quantization grid.
Most semantic variation is then rounded into identical bins.

\paragraph{Online Centering Intervention.}

Instead of waiting for residual normalization,
we insert a zero-centering operator
immediately after mean generation and before quantization:

\[
\mathcal{Z}(Y) = Y - \mathbf{1}\mu_Y^\top.
\]

This enforces

\[
R \approx 1
\]

prior to low-precision casting.

\paragraph{Geometric Benefit.}

Online centering:
\begin{itemize}
\item removes the coherent drift before it accumulates,
\item preserves isotropic semantic variance,
\item maximizes signal-to-quantization-noise ratio (SQNR),
\item prevents cone-shaped collapse in high dimension.
\end{itemize}

This is a surgical intervention:
it removes the structured rank-one component
without altering covariance geometry.

% ============================================================
\subsection{2. Bias-Aware Compilers: Kernel-Level Hardware Alignment}
\label{subsec:bias_aware_compiler}

Naively subtracting the mean introduces additional memory passes,
increasing latency and bandwidth pressure.

To maintain throughput,
bias removal must be fused into the compute kernels themselves.

\paragraph{Fusion-for-Fidelity Strategy.}

Instead of a separate normalization pass,
the compiler modifies the kernel of the nonlinear operator.

Example: SwiGLU kernel fusion.

For each element $y_{ij}$:
\begin{enumerate}
\item accumulate running mean in registers,
\item subtract mean before quantization cast,
\item write centered value directly.
\end{enumerate}

No additional tensor traversal is required.

\paragraph{Preventing Euclidean Drift.}

In low-bit arithmetic,
elementwise error $\epsilon$ accumulates across dimensions:

\[
\|\text{error}\| \sim \sqrt{H}\,\epsilon.
\]

If bias dominates,
quantization scale aligns with $\|\mu\|$,
wasting representational bins tracking shared structure.

Kernel-level centering prevents this amplification
before it propagates across dimensions.

\paragraph{Hardware Efficiency.}

Bias-aware fusion:
\begin{itemize}
\item eliminates extra memory reads/writes,
\item preserves arithmetic intensity,
\item avoids the memory wall bottleneck,
\item scales naturally with tensor cores and SIMD lanes.
\end{itemize}

Thus centering becomes effectively ``free''
from a hardware pipeline perspective.

% ============================================================
\subsection{3. Formal Rationale: The Clean-Block Principle}
\label{subsec:clean_block}

From Theorem~4.1 (Mean Propagation),
bias propagates multiplicatively:

\[
\mathbb{E}[\delta x^\top]
=
\mathbb{E}[\delta]\mu^\top
+
\mathbb{E}[\delta \tilde{x}^\top].
\]

If the block input carries coherent mean,
subsequent linear projections amplify it.

\paragraph{Clean-Block Design.}

We enforce:

\[
\widetilde{X}_\ell = X_\ell - \mu_\ell
\]

before the next linear projection.

Then

\[
\mathbb{E}[\widetilde{X}_\ell] = 0,
\]

and mean amplification in the next layer is prevented.

\paragraph{Dirty Residual Backbone.}

Importantly,
the residual stream need not be erased entirely.
The residual path may preserve accumulated structure,
while the compute block operates on purified inputs.

This produces a dual geometry:
\begin{itemize}
\item residual path: structural carrier,
\item compute block: isotropic semantic processor.
\end{itemize}

\paragraph{Variance Preservation.}

Unlike spike deletion,
centering preserves covariance:

\[
\mathrm{Cov}(X - \mu) = \mathrm{Cov}(X).
\]

Thus semantic geometry remains intact.

% ============================================================
\subsection{4. Physical-World Fidelity Perspective}

For physical modeling tasks,
tiny angular deviations encode forces, conservation laws,
and interaction dynamics.

If every representation is shifted by a dominant bias vector,
the effective angular resolution collapses.

This produces:
\begin{itemize}
\item feature erasure,
\item hallucination-like behavior,
\item loss of fine-grained reasoning.
\end{itemize}

Online centering ensures that
small angular differences remain representable,
even under aggressive quantization.

% ============================================================
\subsection{Summary of Stabilization Principles}

From first principles, stable LLM operator design requires:

\begin{enumerate}
\item Remove coherent mean drift immediately after generation.
\item Prevent residual accumulation of structured bias.
\item Preserve isotropic variance geometry.
\item Align arithmetic conditioning with hardware constraints.
\end{enumerate}

These principles arise directly from:
\begin{itemize}
\item high-dimensional concentration phenomena,
\item nonlinear asymmetry,
\item residual accumulation dynamics,
\item quantization error amplification.
\end{itemize}

Bias mitigation is therefore not merely an optimization trick;
it is a geometric conditioning requirement
for high-dimensional learning under constrained arithmetic.


\section{Conclusion}

Dimensional collapse is not a mysterious emergent pathology.
It is the predictable geometric consequence of:

\begin{itemize}
\item High-dimensional concentration,
\item Nonlinear mean regeneration,
\item Residual accumulation,
\item Finite-precision amplification.
\end{itemize}

Understanding this chain allows architectural and hardware-level mitigation
without destroying semantic information.